\textbf{Phần 1. Trắc nghiệm lựa chọn}
\setcounterpageref{ex}{0}

\begin{ex} %[3H7N43]
Hình tròn có bán kính $r = 6\text{ cm}$ thì đường kính là
\choice
{$6\text{ cm}$}
{$18\text{ cm}$}
{$3\text{ cm}$}
{\True $12\text{ cm}$}
\loigiai{
$6 \times 2 = 12\text{ cm}$.\\
\textbf{Đáp án của câu hỏi là D}
}
\end{ex}

\begin{ex} %[3H7N43]
Khối lập phương có bao nhiêu mặt và bao nhiêu đỉnh?
\choice
{\True $6$ mặt, $8$ đỉnh}
{$8$ mặt, $6$ đỉnh}
{$6$ mặt, $12$ đỉnh}
{$12$ mặt, $8$ đỉnh}
\loigiai{
Khối lập phương có 6 mặt và 8 đỉnh.\\
\textbf{Đáp án của câu hỏi là A}
}
\end{ex}

\begin{ex} %[3H7H43]
Đoạn thẳng $AB = 16\text{ cm}$, $M$ là trung điểm của $AB$. Đoạn $AM$ dài là
\choice
{$6\text{ cm}$}
{\True $8\text{ cm}$}
{$10\text{ cm}$}
{$4\text{ cm}$}
\loigiai{
$16 : 2 = 8\text{ cm}$.\\
\textbf{Đáp án của câu hỏi là B}
}
\end{ex}

\begin{ex} %[3H7H43]
Chu vi hình vuông có cạnh dài $15\text{ cm}$ là
\choice
{$45\text{ cm}$}
{$75\text{ cm}$}
{\True $60\text{ cm}$}
{$30\text{ cm}$}
\loigiai{
$15 \times 4 = 60\text{ cm}$.\\
\textbf{Đáp án của câu hỏi là C}
}
\end{ex}

\begin{ex} %[3H7H43]
Chu vi hình chữ nhật có chiều dài $18\text{ cm}$, chiều rộng $12\text{ cm}$ là
\choice
{$30\text{ cm}$}
{$50\text{ cm}$}
{$40\text{ cm}$}
{\True $60\text{ cm}$}
\loigiai{
$(18 + 12) \times 2 = 60\text{ cm}$.\\
\textbf{Đáp án của câu hỏi là D}
}
\end{ex}

\begin{ex} %[3H7H43]
Đổi: $3\text{ m } 50\text{ mm} = \dots\text{ mm}$.
\choice
{$350\text{ mm}$}
{\True $3050\text{ mm}$}
{$3500\text{ mm}$}
{$3005\text{ mm}$}
\loigiai{
$3000 + 50 = 3050\text{ mm}$.\\
\textbf{Đáp án của câu hỏi là B}
}
\end{ex}

\begin{ex} %[3H7V43]
Một hình tam giác có chu vi $45\text{ cm}$, độ dài hai cạnh là $12\text{ cm}$ và $15\text{ cm}$. Độ dài cạnh thứ ba là
\choice
{\True $18\text{ cm}$}
{$16\text{ cm}$}
{$20\text{ cm}$}
{$22\text{ cm}$}
\loigiai{
$45 - (12 + 15) = 45 - 27 = 18\text{ cm}$.\\
\textbf{Đáp án của câu hỏi là A}
}
\end{ex}

\begin{ex} %[3H7V43]
Có 4 gói bột ngọt, mỗi gói nặng $250\text{ g}$. Tổng khối lượng của 4 gói là
\choice
{$800\text{ g}$}
{$1200\text{ g}$}
{$900\text{ g}$}
{\True $1\text{ kg}$ (hay $1000\text{ g}$)}
\loigiai{
$250 \times 4 = 1000\text{ g} = 1\text{ kg}$.\\
\textbf{Đáp án của câu hỏi là D}
}
\end{ex}

\begin{ex} %[3H7V43]
Một can chứa $2\text{ l}$ nước, người ta rót ra 5 cốc, mỗi cốc $300\text{ ml}$. Lượng nước còn lại trong can là
\choice
{$400\text{ ml}$}
{$600\text{ ml}$}
{\True $500\text{ ml}$}
{$700\text{ ml}$}
\loigiai{
$2000 - 300 \times 5 = 2000 - 1500 = 500\text{ ml}$.\\
\textbf{Đáp án của câu hỏi là C}
}
\end{ex}

\begin{ex} %[3H7V43]
Nhiệt độ buổi sáng là $20^\circ\text{C}$, buổi trưa tăng thêm $9^\circ\text{C}$, chiều tối giảm $4^\circ\text{C}$ so với buổi trưa. Nhiệt độ chiều tối là
\choice
{$24^\circ\text{C}$}
{\True $25^\circ\text{C}$}
{$26^\circ\text{C}$}
{$27^\circ\text{C}$}
\loigiai{
Trưa: $20 + 9 = 29^\circ\text{C}$.\\
Tối: $29 - 4 = 25^\circ\text{C}$.\\
\textbf{Đáp án của câu hỏi là B}
}
\end{ex}

\begin{ex} %[3H7V43]
Một hình vuông có chu vi bằng chu vi hình chữ nhật có chiều dài $14\text{ cm}$, chiều rộng $6\text{ cm}$. Cạnh hình vuông đó là
\choice
{$8\text{ cm}$}
{$12\text{ cm}$}
{\True $10\text{ cm}$}
{$14\text{ cm}$}
\loigiai{
Chu vi: $(14 + 6) \times 2 = 40\text{ cm}$.\\
Cạnh hình vuông: $40 : 4 = 10\text{ cm}$.\\
\textbf{Đáp án của câu hỏi là C}
}
\end{ex}

\begin{ex} %[3H7C43]
Một khung dây sắt hình lập phương có cạnh $15\text{ cm}$. Người ta dùng một đoạn dây sắt dài $2\text{ m}$ để uốn khung đó. Đoạn dây sắt còn thừa bao nhiêu xăng-ti-mét?
\choice
{\True $20\text{ cm}$}
{$15\text{ cm}$}
{$25\text{ cm}$}
{$30\text{ cm}$}
\loigiai{
Khung cần: $15 \times 12 = 180\text{ cm}$.\\
Đoạn dây dài: $2\text{ m} = 200\text{ cm}$.\\
Còn thừa: $200 - 180 = 20\text{ cm}$.\\
\textbf{Đáp án của câu hỏi là A}
}
\end{ex}

\textbf{Phần 2. Trắc nghiệm đúng sai}
\setcounterpageref{ex}{0}

\begin{ex} %[3H7H43]
Xét tính đúng sai về hình học và đại lượng:
\choiceTF
{\True Góc tạo bởi 2 kim đồng hồ lúc 9 giờ đúng là góc vuông.}
{\True $1\text{ kg} = 1000\text{ g}$.}
{Hình tròn đường kính $10\text{ cm}$ có bán kính $20\text{ cm}$.}
{\True Hình tròn đường kính $10\text{ cm}$ có bán kính $5\text{ cm}$.}
\loigiai{
\begin{itemize}
    \item Lúc 9 giờ là góc vuông: Đúng.
    \item $1\text{ kg} = 1000\text{ g}$: Đúng.
    \item Bán kính $r = 10 : 2 = 5\text{ cm} \neq 20\text{ cm}$. Ý 3 Sai, Ý 4 Đúng.
\end{itemize}
\begin{itemize}
    \item a) \textbf{Đúng}.
    \item b) \textbf{Đúng}.
    \item c) \textbf{Sai}.
    \item d) \textbf{Đúng}.
\end{itemize}
}
\end{ex}

\begin{ex} %[3H7V43]
Cho hình chữ nhật có chiều dài $20\text{ cm}$, chiều rộng bằng nửa chiều dài. Xét tính đúng sai:
\choiceTF
{\True Chiều rộng hình chữ nhật là $10\text{ cm}$.}
{\True Chu vi hình chữ nhật là $60\text{ cm}$.}
{Nếu tăng chiều rộng thêm $5\text{ cm}$ thì chu vi tăng thêm $5\text{ cm}$.}
{\True Nếu tăng chiều rộng thêm $5\text{ cm}$ thì chu vi tăng thêm $10\text{ cm}$ ($5 \times 2 = 10$).}
\loigiai{
\begin{itemize}
    \item Rộng: $20 : 2 = 10\text{ cm}$. Đúng.
    \item Chu vi: $(20 + 10) \times 2 = 60\text{ cm}$. Đúng.
    \item Tăng chiều rộng thêm 5 thì chu vi tăng $5 \times 2 = 10\text{ cm}$. Ý 3 Sai, Ý 4 Đúng.
\end{itemize}
\begin{itemize}
    \item a) \textbf{Đúng}.
    \item b) \textbf{Đúng}.
    \item c) \textbf{Sai}.
    \item d) \textbf{Đúng}.
\end{itemize}
}
\end{ex}

\textbf{Phần 3. Trả lời ngắn}
\setcounterpageref{ex}{0}

\begin{ex} %[3H7H43]
Hình tròn có bán kính $9\text{ cm}$. Tính đường kính hình tròn theo xăng-ti-mét.
\shortans{18}
\loigiai{
\textbf{Đáp số:} $18$\\
$9 \times 2 = 18\text{ cm}$.
}
\end{ex}

\begin{ex} %[3H7H43]
Tính chu vi hình vuông có cạnh dài $12\text{ cm}$.
\shortans{48}
\loigiai{
\textbf{Đáp số:} $48$\\
$12 \times 4 = 48\text{ cm}$.
}
\end{ex}

\begin{ex} %[3H7V43]
Tính chu vi hình chữ nhật có chiều dài $16\text{ cm}$, chiều rộng $9\text{ cm}$.
\shortans{50}
\loigiai{
\textbf{Đáp số:} $50$\\
$(16 + 9) \times 2 = 50\text{ cm}$.
}
\end{ex}

\begin{ex} %[3H7V43]
Đoạn thẳng $CD = 30\text{ cm}$, $M$ là trung điểm của $CD$. Tính độ dài đoạn $CM$ theo xăng-ti-mét.
\shortans{15}
\loigiai{
\textbf{Đáp số:} $15$\\
$30 : 2 = 15\text{ cm}$.
}
\end{ex}

\begin{ex} %[3H7V43]
Đổi: $5\text{ cm } 6\text{ mm}$ ra mi-li-mét.
\shortans{56}
\loigiai{
\textbf{Đáp số:} $56$\\
$50 + 6 = 56\text{ mm}$.
}
\end{ex}

\begin{ex} %[3H7C43]
Một hình chữ nhật có chu vi $36\text{ cm}$, chiều dài gấp đôi chiều rộng. Tính chiều dài hình chữ nhật theo xăng-ti-mét.
\shortans{12}
\loigiai{
\textbf{Đáp số:} $12$\\
Nửa chu vi: $36 : 2 = 18\text{ cm}$.\\
Tổng số phần: $1 + 2 = 3$ phần.\\
Chiều dài: $18 : 3 \times 2 = 12\text{ cm}$.
}
\end{ex}
